% ============================================================
%  AUTONOMY, CAPACITY, AND THE RIGHT TO EXIT
%  Master Policy Packet — Non-Terminal Medical Aid in Dying
%  Justin Bogner, 2026
%  Compile: xelatex 00_master_policy.tex (twice for TOC/refs)
% ============================================================

\documentclass[11pt, letterpaper]{article}

% ── Core ────────────────────────────────────────────────────
\usepackage{fontspec}
\usepackage{microtype}
\usepackage[margin=1.15in, bottom=1.2in, top=1.1in, headheight=22pt]{geometry}
\usepackage{setspace}
\onehalfspacing

% ── Fonts ───────────────────────────────────────────────────
\setmainfont{PT Serif}
\setsansfont{PT Sans}

% ── Structure ───────────────────────────────────────────────
\usepackage{titlesec}
\usepackage{titletoc}
\usepackage{fancyhdr}
\usepackage{booktabs}
\usepackage{tabularx}
\usepackage{threeparttable}
\usepackage{enumitem}
\usepackage{array}
\usepackage{longtable}
\usepackage{multirow}

% ── Visual ──────────────────────────────────────────────────
\usepackage{xcolor}
\usepackage{mdframed}
\usepackage{tcolorbox}
\tcbuselibrary{skins, breakable}
\usepackage{graphicx}

% ── Links ───────────────────────────────────────────────────
\usepackage[hidelinks, colorlinks=false, pdfnewwindow=true]{hyperref}

% ── Footnotes ───────────────────────────────────────────────
\usepackage[hang, flushmargin]{footmisc}
\setlength{\footnotemargin}{1.2em}

% ─────────────────────────────────────────────────────────────
%  COLOR PALETTE
% ─────────────────────────────────────────────────────────────
\definecolor{accent}{HTML}{1B3A6B}
\definecolor{accentlight}{HTML}{2E5D9E}
\definecolor{calloutbg}{HTML}{EAF0FB}
\definecolor{calloutborder}{HTML}{2E5D9E}
\definecolor{warningbg}{HTML}{FDF3E7}
\definecolor{warningborder}{HTML}{C0622F}
\definecolor{darkgray}{HTML}{2B2B2B}
\definecolor{midgray}{HTML}{555555}
\definecolor{rulegray}{HTML}{BBBBBB}

% ─────────────────────────────────────────────────────────────
%  SECTION FORMATTING
% ─────────────────────────────────────────────────────────────
\renewcommand{\thesection}{\Roman{section}}
\renewcommand{\thesubsection}{\arabic{section}.\arabic{subsection}}

\titleformat{\section}
  {\sffamily\large\bfseries\color{accent}\centering}
  {\thesection.}{0.6em}{}[{\color{rulegray}\titlerule[0.5pt]}]

\titleformat{\subsection}
  {\sffamily\normalsize\bfseries\color{darkgray}}
  {\thesubsection}{0.5em}{}

\titlespacing*{\section}{0pt}{2.2ex plus 1ex minus 0.2ex}{1.2ex plus 0.2ex}
\titlespacing*{\subsection}{0pt}{1.8ex plus 0.8ex}{0.6ex plus 0.1ex}

% ─────────────────────────────────────────────────────────────
%  PAGE STYLE
% ─────────────────────────────────────────────────────────────
\pagestyle{fancy}
\fancyhf{}
\fancyhead[L]{\small\sffamily\color{midgray} Autonomy, Capacity, and the Right to Exit}
\fancyhead[R]{\small\sffamily\color{midgray} Master Policy Packet}
\fancyfoot[C]{\small\sffamily\color{midgray}\thepage}
\renewcommand{\headrulewidth}{0.4pt}
\renewcommand{\headrule}{\hbox to\headwidth{\color{rulegray}\leaders\hrule height \headrulewidth\hfill}}
\renewcommand{\footrulewidth}{0pt}

% ─────────────────────────────────────────────────────────────
%  CALLOUT ENVIRONMENTS
% ─────────────────────────────────────────────────────────────
\tcbset{
  policybox/.style={
    enhanced, breakable,
    colback=calloutbg, colframe=calloutborder,
    boxrule=0.5pt, leftrule=4pt,
    left=10pt, right=10pt, top=8pt, bottom=8pt,
    arc=0pt, outer arc=0pt,
    fontupper=\small,
  },
  warningbox/.style={
    enhanced, breakable,
    colback=warningbg, colframe=warningborder,
    boxrule=0.5pt, leftrule=4pt,
    left=10pt, right=10pt, top=8pt, bottom=8pt,
    arc=0pt, outer arc=0pt,
    fontupper=\small,
  },
  statutebox/.style={
    enhanced, breakable,
    colback=white, colframe=accent,
    boxrule=0.8pt,
    left=12pt, right=12pt, top=10pt, bottom=10pt,
    arc=0pt, outer arc=0pt,
    fontupper=\small\ttfamily,
  }
}

% ─────────────────────────────────────────────────────────────
%  LIST STYLES
% ─────────────────────────────────────────────────────────────
\setlist[itemize,1]{
  leftmargin=1.4em, label={\color{accentlight}\textbullet},
  itemsep=3pt, parsep=0pt, topsep=4pt
}
\setlist[enumerate,1]{
  leftmargin=1.6em,
  itemsep=4pt, parsep=0pt, topsep=4pt
}

% ─────────────────────────────────────────────────────────────
%  METADATA
% ─────────────────────────────────────────────────────────────
\title{%
  {\sffamily\LARGE\bfseries\color{accent}
    AUTONOMY, CAPACITY, AND THE RIGHT TO EXIT}\\[0.6em]
  {\large\itshape A Policy Framework for Non-Terminal Medical Aid in Dying}\\[0.4em]
  {\normalsize\sffamily\color{midgray} Master Policy Packet — For Distribution to Legislative,
    Clinical, Academic, and Advocacy Audiences}
}
\author{\sffamily\bfseries Justin Bogner}
\date{\sffamily 2026}

% ============================================================
\begin{document}
% ============================================================

\maketitle
\thispagestyle{empty}

\vspace{0.5em}
{\color{rulegray}\hrule height 0.5pt}
\vspace{1.2em}

% ── AUTHOR'S NOTE ────────────────────────────────────────────
\begin{tcolorbox}[policybox, title={\sffamily\bfseries\small\color{accent} Author's Note}]
This framework emerges from lived contradiction: at eighteen, I survived a serious suicide
attempt with full intent and means. Nearly twenty years later, I choose life---not out of
fear or social obligation, but because the conditions became tolerable.

\medskip
That trajectory does not invalidate the argument; it exposes the lie that forced preservation
is benevolence. The state that resuscitates without consent does not save a life; it claims
ownership of one. This paper insists the right to exit belongs to the person whose life it
is. Only when the door stands unlocked does the choice to remain carry moral weight. Anything
else is conscription.

\medskip
\textit{This packet is the source document for four audience-specific briefs: \href{01_legislative_brief.pdf}{Legislative},
\href{04_clinical_brief.pdf}{Clinical/Medical Ethics}, \href{03_academic_paper.pdf}{Academic}, and \href{02_advocacy_brief.pdf}{Advocacy}. All derive from the argument presented here.}
\end{tcolorbox}

\vspace{1em}

% ── TABLE OF CONTENTS ────────────────────────────────────────
{
\hypersetup{hidelinks}
\tableofcontents
}

\newpage

% ============================================================
\section*{Executive Summary}
\addcontentsline{toc}{section}{Executive Summary}
% ============================================================

Current Medical Aid in Dying (MAiD) frameworks predicate eligibility on terminal illness,
irremediable suffering, or analogous threshold conditions. This paper argues that these
criteria are philosophically insufficient and practically counterproductive. The operative
criterion for eligibility should be \textbf{decisional capacity}---and nothing more.

\medskip
This framework proposes a shift from a \textbf{Protectionist Model}---which treats eligibility
as a concession granted when suffering becomes sufficiently legible to clinicians and
legislators---to a \textbf{sovereignty model} grounded in four claims:

\begin{enumerate}
  \item \textbf{Capacity is the only gate.} Suffering is a proxy; capacity is the
    principle. A person without pain who has formed a settled, informed preference for
    death presents an equal autonomous claim to one in persistent agony. The state has
    no legitimate interest in ranking whose suffering is sufficiently legible.

  \item \textbf{Process serves as the filter.} A rigorous two-stage assessment and
    mandatory waiting period distinguishes crisis response from stable autonomous
    preference. Procedure is not a barrier; it is the mechanism of verification.

  \item \textbf{Capability does not create obligation.} An individual's value to
    family, community, or society does not function as a statutory prohibition on exit.
    Persons are not public assets. Continued existence is not a debt owed to the collective.

  \item \textbf{The Right to Exit as forcing function.} An enforceable exit right
    compels the state to treat livability as a deliverable rather than a platitude.
    Where socioeconomic conditions drive MAiD uptake, those conditions become documented
    public record. The moral and political cost of that record creates structural
    incentive for investment in life. This is not advocacy for death; it is the most
    consequential structural argument for dignified living.
\end{enumerate}

\begin{tcolorbox}[warningbox]
\textbf{On the disability rights objection:} This framework directly engages, rather than
dismisses, the critique advanced by organizations such as DAWN Canada and ADAPT. The
principal concern---that expanding MAiD eligibility enables the state to offer death as a
substitute for adequate supports---is treated as \emph{partially valid and structurally
important}. The forcing function argument presented in Section~V is offered as the
substantive answer to that concern, not a circumvention of it. See Section~IV for full
engagement.
\end{tcolorbox}

\newpage

% ============================================================
\section{The Problem with Current Frameworks}
% ============================================================

\subsection{The Suffering Criterion}

Jurisdictions that permit MAiD typically require proof of ``unbearable'' or ``intolerable''
suffering. This framing conflates sympathy with rights. It suggests that a person must first
be visibly broken before they are permitted to leave.

The suffering criterion creates a perverse dynamic: the more articulately a person can describe
their distress to an evaluator, the more likely they are to receive access. Conversely, a
person who has formed a settled, informed, philosophically coherent preference for death---but
who presents as calm and functional---may be denied on the grounds that their suffering is
insufficiently demonstrable. This is not a clinical standard; it is a performance requirement.

A more fundamental problem is definitional. ``Unbearable'' is not a clinical category. It is
a subjective threshold that varies with individual psychology, cultural background, pain
tolerance, and access to support services. To adjudicate on it requires evaluators to
override the applicant's own report in favor of external assessment---precisely the
paternalism this framework rejects.

\subsection{The Terminal Illness Requirement}

The United States model, operative in ten states and the District of Columbia, restricts
eligibility to persons with a terminal prognosis of six months or fewer. This restriction
has no coherent philosophical foundation.

The relevant moral question is not \emph{how soon} a person will die, but whether they possess
the autonomous capacity to determine when and how. A person with a terminal prognosis of
five months is not meaningfully different, in terms of self-determination, from a person with
a non-terminal but severely degenerative condition who will live twenty years in declining
function. To treat them differently is to make a rights determination on the basis of
prognosis rather than personhood.

\subsection{International Comparators}

\begin{table}[h]
\centering
\small
\begin{threeparttable}
\begin{tabularx}{\textwidth}{l X X X}
\toprule
\textbf{Jurisdiction} & \textbf{Eligibility} & \textbf{Suffering Req.?} & \textbf{Non-Terminal?} \\
\midrule
Canada (2026)
  & Grievous \& irremediable condition; Track~2 for non-foreseeable death
  & Yes
  & Yes (Track~2; 732 provisions in 2024; 61.5\% disability)$^{\dagger}$ \\[4pt]
United States
  & Terminal ($\leq$6 months); 10 states + DC
  & No
  & No \\[4pt]
Belgium / Netherlands
  & Unbearable suffering including psychiatric
  & Yes
  & Yes (psychiatric review) \\[4pt]
Switzerland
  & Any (assisted suicide organizations)
  & No
  & Yes (existential cases) \\[4pt]
\textbf{Proposed}
  & \textbf{Decisional capacity only}
  & \textbf{No}
  & \textbf{Yes (stable preference)} \\
\bottomrule
\end{tabularx}
\begin{tablenotes}
  \footnotesize
  \item[$\dagger$]
    Health Canada, \textit{Sixth Annual Report on Medical Assistance in Dying}
    (2024 data, released November 2025). Mental illness as sole underlying condition
    delayed to March 2027. Track~2 year-over-year growth: 17\% from 625 (2023) to 732 (2024).
\end{tablenotes}
\end{threeparttable}
\caption{MAiD Eligibility Models, 2026}
\label{tab:comparators}
\end{table}

\subsection{Liability Theater}

Psychiatric response to suicidal ideation frequently functions as institutional liability
management rather than therapeutic care. The operational priority is avoiding documented
responsibility for a death ``on our watch,'' not resolving the conditions that produced the
crisis. The result is a sequence that is coercive in substance if not in name: involuntary or
pressured hospitalization, coerced safety contracts of no demonstrated clinical value,
premature discharge into the same unlivable conditions, and the restarting of the cycle.

Biology persists; autonomy is suspended. The patient is held as a hostage to the provider's
liability exposure and the state's moral discomfort with the existence of the preference.
This is not a treatment model. It is an administrative loop designed to defer, not resolve.

\newpage

% ============================================================
\section{The Philosophical Foundation}
% ============================================================

\subsection{Bodily Self-Determination}

Medical jurisprudence across Western democracies has established, without serious remaining
controversy, that competent adults may refuse any medical treatment---including life-sustaining
treatment---without requiring justification, without being required to attempt it first, and
without the consent of family members or treating clinicians. The foundation of this right
is bodily self-determination: the principle that a person's physical existence is their own
domain, not the state's.

If this right is recognized---and it is, in law, in every relevant jurisdiction---then the
question is not whether persons possess the right to determine the boundaries of their
existence, but why that right is suspended the moment the determination is self-directed
rather than treatment-directed. The moral logic that permits a patient to refuse a ventilator
and die by suffocation also supports permitting that patient to choose a swifter and less
distressing death by other means. The distinction the law currently draws is not principled;
it is aesthetic.

\subsection{The Informed Refusal Principle}

Medicine already recognizes, without controversy, that a patient may decline any
treatment---including life-sustaining treatment---without first attempting it. A patient may
refuse chemotherapy without a single round. A patient may decline dialysis without trial.
The sole requirement is that the refusal be \textit{informed}: the person must understand
what the treatment is, what outcomes it realistically offers, and what declining means.

This framework applies that principle consistently and without exception.

An individual seeking MAiD assessment need not have attempted psychiatric treatment, exhausted
pharmacological options, or completed a course of psychotherapy. They must know that these
interventions exist, understand their realistic outcomes and risks, and---knowing this---decline
them. The criterion is \textbf{informed awareness}, not \textbf{attempted compliance}.

\begin{tcolorbox}[policybox]
\textbf{The clinical interview's role under this framework:} The evaluating clinician tests
comprehension, not compliance history. The relevant questions are not ``Have you tried an
SSRI?'' but ``What do you understand about SSRIs, what outcomes do the literature suggest,
and what led you to decline them?'' An honest, documented refusal following full disclosure
is more epistemically transparent than a coerced trial producing performative compliance.
\end{tcolorbox}

\medskip
Requiring prior exhaustion of treatment for psychiatric or existential cases imposes a double
standard that medicine has abandoned everywhere else. It encodes a covert hierarchy: the
state's preferred interventions become mandatory prerequisites for autonomy in the categories
it deems ``weak-minded.'' This is not diligence. It is diagnostic disenfranchisement.

\subsection{Capability Does Not Create Obligation}

A persistent objection is that an individual's value to others---family, society, their
professional community---prohibits their exit. This view is structurally identical to the
claim that a person's labor value prohibits their resignation. Both treat the person as a
resource to be retained by others' need rather than a subject with sovereign interest in
their own existence.

Relationships are vital. They are, for many people, primary reasons to remain. But they cannot
function as statutory chains. The fact that a person \textit{can} contribute does not mean
they \textit{must}. Continued existence is not a debt owed to the collective, to one's family,
or to the state.

\subsection{The Depressive Capacity Objection}

A targeted form of the paternalist objection holds that persons with treatment-resistant
depression cannot satisfy the ``appreciation of alternatives'' criterion for capacity,
because hopelessness functions as a cognitive filter that makes alternatives appear
unavailable even when they are formally understood.

This objection deserves direct engagement. The informed refusal framework dissolves it by
clarifying what appreciation requires. The question is not whether the person \emph{feels}
the alternatives to be viable---this would make the capacity standard dependent on the
very symptom being assessed---but whether they \emph{understand} what the alternatives are
and what they offer.

A preference that survives full informed disclosure is the operationalization of capacity.
To permanently exclude persons with depressive histories on the grounds that their refusals
cannot be trusted is to convert a diagnosis into a lifetime disqualifier. This is the
paternalism this framework rejects, made explicit.

\newpage

% ============================================================
\section{The Two-Stage Capacity Assessment}
% ============================================================

To distinguish authentic autonomous preference from crisis-state decision-making, this
framework proposes a two-stage protocol with a mandatory assessment period.

\subsection{Stage 1: Decisional Capacity}

The evaluating clinician must assess and document that the applicant demonstrates:

\begin{itemize}
  \item \textbf{Understanding of consequences.} The applicant understands the permanent and
    irreversible nature of death. This must be assessed, not assumed.

  \item \textbf{Informed awareness of alternatives.} The applicant has received documented
    presentation of all available alternatives---medical, psychiatric, financial, social,
    palliative---and understands their realistic outcomes and associated risks. Prior
    utilization of any alternative is \emph{not} required. Informed refusal of alternatives,
    following full disclosure, must be respected.

  \item \textbf{Absence of acute distortion.} The applicant is not in an active psychotic
    state or transient crisis condition that materially undermines comprehension. This is a
    state assessment, not a diagnostic exclusion: a person with a history of psychosis who
    is currently stable is not disqualified.

  \item \textbf{Preference stability.} The preference is stable across time and is not
    reducible to a situational or episodic response to a transient condition. This is assessed
    through the waiting period and multiple evaluations, not through a single interview.
\end{itemize}

\subsection{Stage 2: Process Comprehension}

The applicant must demonstrate willingness to engage with the assessment process itself,
including understanding \emph{why} the safeguards exist and accepting them as legitimate.
This functions as a mechanical filter for impulsivity.

A person in acute crisis will rarely tolerate a 90-day bureaucratic process with
multiple independent evaluations. A person with a resolved, stable autonomous preference can.
The process itself is evidence of resolution---not because endurance proves rationality, but
because the capacity to engage with a structured procedure across time reflects precisely
the stable, reflective orientation that distinguishes autonomous preference from reactive crisis.

\begin{tcolorbox}[policybox]
\textbf{Diagnostic logic of Stage 2:} A person who claims Stage~1 capacity while refusing
to engage with Stage~2 process requirements likely fails Stage~1. Genuine understanding of
why safeguards exist---as protections for the applicant, not obstacles---is itself a
component of the comprehension required for autonomous decision-making.
\end{tcolorbox}

\newpage

% ============================================================
\section{Principal Objections and Responses}
% ============================================================

\subsection{The Disability Rights Critique}

\begin{tcolorbox}[warningbox, title={\sffamily\bfseries\small\color{warningborder}
  Objection — DAWN Canada, ADAPT, and related organizations}]
Expanding MAiD eligibility, particularly to non-terminal and socioeconomically disadvantaged
applicants, does not extend autonomy---it provides the state with a cheaper alternative to
adequate supports. Disabled people, people experiencing poverty, and people who are isolated
are not choosing death freely; they are choosing death because the state has failed to make
life livable. An expanded eligibility framework legitimizes that failure by processing it
as autonomous preference.
\end{tcolorbox}

\medskip
\noindent\textbf{Response: The critique is partially correct, and that is precisely why this
framework does not stop at expanding eligibility.}

The disability rights organizations advancing this objection have identified a real phenomenon.
Health Canada's 2024 data shows that Track~2 recipients are modestly more likely to reside
in lower-income neighborhoods than Track~1 recipients, and that disability and social isolation
are prominent documented drivers of ``intolerable suffering'' claims. The Ontario Chief
Coroner's reviews have documented cases in which inadequate housing, income insufficiency, and
care gaps appear as primary conditions motivating MAiD requests.

This is not acceptable. The question is what follows from that finding.

The Protectionist position---restrict access, close the exit door---does not resolve the
underlying conditions. It adds legal coercion to material deprivation. The person who lacks
housing, adequate care, and social connection, and who finds life intolerable as a result,
is not made better off by being denied the right to act on that assessment. They are made
worse off in a way that is now compelled by the state.

The forcing function argument in Section~V of this paper is the substantive answer to the
disability rights concern, not a deflection of it:

\begin{enumerate}
  \item Under the proposed framework, evaluators are \emph{required} to present all available
    alternatives---housing, income support, disability services, community connection---with
    sufficient specificity for genuine comprehension. This presentation must be documented.
  \item Where an applicant cites inadequate access to these supports as a condition of their
    preference, the administering authority is \emph{required} to transmit a formal report
    to relevant government ministries documenting this as a driver of the provision.
  \item That documentation creates a public record of state failure. Not a private administrative
    note---a documentable, aggregable, politically legible record.
\end{enumerate}

The disability rights organizations are right that the state should not offer death as an
alternative to adequate care. This framework makes that explicit and structural. Where they
are wrong is in the conclusion: that the remedy is to remove the exit rather than to document,
publicize, and politically weaponize the conditions that drive people toward it.

The Protectionist Model has governed most jurisdictions for decades. In those same decades,
housing supports have eroded, disability services have been underfunded, psychiatric wait
times have grown. Closing the exit has not produced the dignified conditions its defenders
claim it would compel. There is no evidence that coerced existence creates political pressure
for better lives. There is evidence that documented, state-recorded deaths attributable to
inadequate supports create that pressure.

\subsection{The Slippery Slope Objection}

The concern is that expanding eligibility criteria leads inevitably to expansion of the
eligible population in ways that cannot be controlled---that ``capacity only'' becomes
``effectively anyone.''

This objection is empirically testable. The relevant data comes from Belgium and the
Netherlands, which have operated under broad eligibility criteria---including psychiatric
cases---for over two decades. Neither jurisdiction has experienced the predicted exponential
expansion or the predicted collapse of safeguards. Utilization has grown modestly and
measurably; the demographics of recipients have shifted somewhat; the procedural requirements
have been tightened in response to specific case-by-case concerns.

The slippery slope argument proves too much. Applied consistently, it would prohibit any
extension of any right on the grounds that extension is logically unlimited. The answer is
not to prohibit extension but to design the procedural architecture carefully---which this
framework does.

\subsection{The External Coercion Objection}

The concern is that persons---particularly disabled people, elderly people, and people who
experience themselves as burdens---may choose MAiD under subtle or explicit coercive pressure
from family members, caregivers, or financial incentives.

This is a legitimate concern and is addressed procedurally rather than by restricting
eligibility:

\begin{itemize}
  \item Every assessment includes a confidential interview with a social worker or independent
    advocate unaffiliated with the care team.
  \item Financial and dependency circumstances are reviewed as a standard component.
  \item Mandatory offer of independent legal assistance for advance planning.
  \item Revocation is absolute and available until the final moment.
\end{itemize}

The coercion concern is an argument for robust process, not for categorical exclusion.
Categorical exclusion on the basis that coercion \emph{might} be present is itself a coercive
act: it overrides the expressed preference of the many to protect against the hypothetical
victimization of the few.

\newpage

% ============================================================
\section{Systemic Consequences: The Right to Exit as Forcing Function}
% ============================================================

The most profound implication of this framework is not the death of the individual but the
reformation of the society that remains.

\subsection{From Coercion to Alliance}

In a system where the patient possesses the legal right to exit, the clinician loses the
power of coercive retention. They cannot function as a custodial jailer mandating existence.
Instead, the clinician must become an ally---working to make the patient's life subjectively
worth continuing. Therapy shifts from preventing death as an administrative outcome to
incentivizing life as a collaborative project. Custodial becomes restorative. This is not
a small shift; it is a complete reorientation of the therapeutic relationship.

\subsection{The Mandate for Livability}

Health Canada's 2024 data already shows the dynamic in embryonic form: Track~2 provisions
reached 732 in 2024, up 17\% year-over-year, with 61.5\% of recipients reporting disability.
Total MAiD provisions reached 16,499 in 2024, representing 5.1\% of all Canadian deaths.
Documented cases cite unaffordable housing, inadequate disability supports, and social
isolation as core drivers. These are not anomalies. They are early signals.

Under capacity-based eligibility, the signal becomes undeniable. The state must either
document systemic abandonment as a cause of death or prevent it. If a citizen chooses MAiD
because they lack housing or adequate care, the administering authority records that failure.
Aggregated across a population, that record is a political instrument of the first order.

\begin{tcolorbox}[policybox]
\textbf{The Survivor's Dividend:} Resources reallocated to ensure that the choice to remain
is supported by dignity, not merely permitted by default. The Right to Exit functions as a
consumer protection standard for existence: the product must be good enough to retain the
user, or the state must explain, on public record, why it was not.
\end{tcolorbox}

\subsection{Correction by Consent}

It must be acknowledged directly: this framework will result in what may be termed a
\textit{correction by consent}. Some individuals will choose to exit rather than wait for
societal improvement that may not arrive in time for them. This is the price of a free
society---and the price of having failed them earlier.

The alternative is to compel their continued existence while continuing to fail them. That
is not a moral position. It is a policy of managed despair with administrative gatekeeping.

A society that remains by consent rather than by conscription is not a depleted society; it
is a more honest one. And it is a society with an undeniable structural incentive to improve
the conditions under which the consent to remain is given.

\newpage

% ============================================================
\section{Procedural Architecture}
% ============================================================

\subsection{Mandatory Resource Presentation}

Evaluators must present, in sufficient specificity for genuine comprehension, every realistic
alternative: housing assistance, income support, disability services, community connection,
palliative care, psychiatric intervention, peer support programs. Presentation must be
documented. Referral must be offered and, if declined, that declination must be recorded with
the applicant's stated reasons.

Informed refusal of any alternative, once its nature and realistic outcomes are understood,
must be respected without condition.

\subsection{Waiting Periods as Diagnostic Instruments}

A minimum assessment period of 90 days applies to all non-terminal applications. This period
is not a delay tactic; it is a diagnostic instrument. A preference that cannot survive 90 days
of reflection, multiple independent evaluations, and genuine engagement with alternatives is
not the settled autonomous preference this framework is designed to recognize.

The period serves a second function: it is the interval during which alternative supports
may be arranged and tried. Persons who accept alternative support during this period may
withdraw the application without prejudice or stigma at any point.

\subsection{Independent Coercion Screening}

Every assessment includes:

\begin{itemize}
  \item Confidential interview by a social worker or independent patient advocate unaffiliated
    with the treating team and with no institutional stake in the outcome.
  \item Structured review of financial circumstances, dependency relationships, and any
    indicators of external pressure.
  \item Mandatory offer of independent legal assistance for advance planning: wills, powers
    of attorney, care directives.
\end{itemize}

\subsection{Documentation of Socioeconomic Drivers}

Where an applicant identifies inadequate access to housing, income, disability supports,
or social connection as a primary or contributing condition of their preference, the
administering authority \textbf{must} transmit a confidential formal report to the relevant
government ministry, documenting these conditions as a driver of the provision. This is not
optional. This is the mechanism by which the forcing function operates.

\subsection{Revocation}

The right to revoke the application is absolute. It may be exercised at any point in the
process, for any reason, without explanation, without stigma, and without effect on
future access to services. No ``death day locked-in'' mechanism is permissible.

\newpage

% ============================================================
\section{Model Statutory Language}
% ============================================================

The following provisions are offered as draft model language for legislative consideration.
They are not jurisdiction-specific and require adaptation to local constitutional and
statutory frameworks.

\medskip

\begin{longtable}{>{\bfseries\sffamily\small}p{3cm} p{10.5cm}}
\toprule
\textbf{\sffamily Provision} & \textbf{\sffamily Proposed Language} \\
\midrule
\endhead
\bottomrule
\endfoot

Eligible Person
  & Any adult who demonstrates decisional capacity as assessed under this Act.
    Terminal diagnosis and proof of suffering are not required. \\[6pt]

decisional capacity
  & The demonstrated ability to: (a)~understand the permanent and irreversible nature
    of death; (b)~demonstrate informed awareness of available alternatives including
    medical, psychiatric, financial, and social supports; (c)~express a preference
    stable across the assessment period; and (d)~be free from acute psychotic or
    crisis states that materially distort comprehension.
    A history of psychiatric diagnosis does not constitute disqualification.
    Stability is assessed at the time of evaluation. \\[6pt]

Informed Refusal
  & An applicant need not have utilized any alternative prior to application.
    Documented informed awareness of alternatives and reasons for declining
    them shall satisfy this element. Evaluators shall assess comprehension
    of alternatives, not compliance history. \\[6pt]

Assessment Period
  & A minimum of 90 days from first application for non-terminal provisions.
    Multiple independent evaluations are required within this period.
    Applications may be withdrawn at any point without prejudice. \\[6pt]

Coercion Screening
  & Every application shall include: (a)~confidential assessment by an
    independent social worker or patient advocate unaffiliated with the
    care team; (b)~structured review of financial and dependency
    circumstances; (c)~mandatory offer of independent legal assistance
    for advance planning. \\[6pt]

Socioeconomic Documentation
  & Where an applicant identifies inadequate housing, income, disability
    support, or social connection as a primary or contributing basis for
    their preference, the administering authority shall transmit a formal
    report to the relevant ministry documenting these conditions as a
    driver of the provision. \\[6pt]

Revocation
  & The right to revoke is absolute, available without condition or
    explanation until the moment of provision, and has no adverse effect
    on future access to services or applications. \\[6pt]

Clinician Conscience
  & No clinician is required to participate in MAiD assessment or provision.
    Conscientious objectors must provide timely referral to a willing
    provider. Refusal to refer constitutes patient abandonment. \\[6pt]

Standard of Proof
  & Decisional capacity must be established by clear and convincing evidence
    across the full assessment period, consistent with the standard applied
    in guardianship and high-stakes competency determinations.
    The burden rests on the assessment record. \\[6pt]

Evaluator Disagreement
  & Where independent evaluators reach divergent conclusions on capacity,
    the matter is referred to an independent Clinical Ethics Committee
    within fourteen days. Capacity requires reasoned concurrence of at
    least two evaluators. Unresolved disagreement triggers administrative
    review. \\[6pt]

Review and Appeal
  & Any applicant whose request is denied has the right to: (a)~a fresh
    independent evaluation by a new panel within 30 days; (b)~expedited
    administrative tribunal review; and (c)~judicial review of any adverse
    tribunal decision. All denials must be in writing with reasoned findings.
    The administering authority bears the burden of justification on review. \\

\end{longtable}

% ============================================================
\section*{Conclusion: Life as an Opt-In System}
\addcontentsline{toc}{section}{Conclusion}
% ============================================================

This is not advocacy for death.

It is a refusal to treat adult sovereignty as conditional on suffering thresholds, clinical
legibility, or social utility. The argument does not require that exit be easy; it requires
that exit be possible---and that the conditions under which people choose to remain are
documented, legible, and subject to political accountability.

Bar the exit, and presence is conscription: silent, coerced, stripped of moral weight. Unlock
it, and presence becomes deliberate affirmation. Every person who remains does so as a
choice, not a default. That is the only condition under which the choice means anything.

The alternative is gentler rhetoric concealing a rising body count that we already refuse to
name honestly. Replace sanctity-of-life dogma with sovereignty of the living, or acknowledge
that the dogma was always about the comfort of those who remain rather than the welfare of
those who suffer.

The door must stand unlocked. What happens next---individually, socially, politically---is
then a matter of genuine choice. That is where moral life begins.

\newpage

% ============================================================
\section*{References}
\addcontentsline{toc}{section}{References}
% ============================================================

\begin{flushleft}
\setlength{\parindent}{0pt}
\setlength{\parskip}{6pt}

Health Canada. (2025). \textit{Sixth Annual Report on Medical Assistance in Dying in
Canada} (2024 data). Ottawa: Health Canada.

\medskip
Health Canada. (2023--2024). \textit{Track~2 case documentation and socioeconomic
analyses}. Ottawa: Health Canada.

\medskip
Ontario Office of the Chief Coroner. (2022--2025). \textit{Reviews of MAiD cases
involving socioeconomic drivers}. Toronto.

\medskip
United Nations Special Rapporteur on the Rights of Persons with Disabilities. (2019).
\textit{Report on autonomy, legal capacity, and MAiD}. Geneva: UN Human Rights Council.

\medskip
Government of Canada. (2024). \textit{Bill C-7: An Act to amend the Criminal Code
(medical assistance in dying)}. Ottawa: Parliament of Canada.

\medskip
Donnelly, M., \& Campbell, A. (2024). Decisional capacity and the limits of
paternalism in assisted dying law. \textit{Journal of Medical Ethics}, 50(3), 182--191.

\medskip
Disability Action Working Network (DAWN) Canada. (2023). \textit{MAiD and the
inadequacy of social supports: A disability rights analysis}. Ottawa: DAWN Canada.

\medskip
Beauchamp, T.~L., \& Childress, J.~F. (2019). \textit{Principles of Biomedical
Ethics} (8th ed.). Oxford University Press.

\medskip
Carter v.\ Canada (Attorney General), 2015 SCC 5. Supreme Court of Canada.

\end{flushleft}

% ============================================================
\end{document}
% ============================================================
